\begin{definition} \label{definition:tate_twist_of_a_relative_chow_motive_over_a_noetherian_base_and_commutative_ring}
Let $S$ be a \CrefAndHyperrefIfExist{definition:noetherian_scheme}{Noetherian base scheme} and $R$ a \CrefAndHyperrefIfExist{definition:commutative_ring}{commutative ring}. 
In the \CrefAndHyperrefIfExist{definition:category_of_relative_chow_motives_over_a_noetherian_base_and_coefficients_in_a_commutative_ring}{category of relative Chow motives $\mathfrak{Chow}(S, R)$}, the \hldef{Tate twist} is an operational shift denoted by \hl{$(n)$} for $n \in \mathbb{Z}$. For a \CrefAndHyperrefIfExist{definition:relative_chow_motive_over_a_noetherian_base_and_commutative_ring_of_coefficients}{relative Chow motive} $M = (X, p, i)$, the twist is defined as $M(n) = (X, p, i+n)$. Geometrically, the motive of the \CrefAndHyperrefIfExist{definition:projective_space_of_dimension_n_over_a_scheme}{projective line} $\mathbb{P}^1_S$ decomposes as $h(\mathbb{P}^1_S) \cong \mathbb{1}_S \oplus \mathbb{1}_S(1)$, where $\mathbb{1}_S(1)$ is the Lefschetz motive.
\end{definition}